%! Introduction to Polynomial Functions — Unit 2, Lesson 2.0 — Guided Notes (Answer Key)
\documentclass[10pt]{article}
\usepackage{precalculus-article}
\usepackage{precalculus-key}   % pulls in -boxes; do NOT also load -boxes
\newcommand{\TallMath}[1]{$\displaystyle #1\rule[-1.4em]{0pt}{3.2em}$}

\begin{document}

\pageheader{Unit 2, Lesson 2.0}{Guided Notes --- Answer Key}

\vspace{0.12in}

\begin{objectivebox}
By the end of this lesson, I will be able to\ldots
\begin{enumerate}[itemsep=2pt]
  \item Decide whether a given expression defines a \ans{polynomial} function, and
        say why an expression that fails does fail.
  \item Name the \ans{degree}, leading term, leading \ans{coefficient}, and constant
        term of a polynomial, and write one in \ans{standard} form.
  \item Name the low-degree families and count \ans{turning} points on a
        pre-drawn polynomial graph.
\end{enumerate}
\end{objectivebox}

\vspace{0.1in}

\boxguard
\begin{vocabbox}
Fill in each term as we define it together.
\par\vspace{2pt}
\noindent\textbf{\textcolor{plum}{Polynomial function:}}\\[1pt]
\ansline{A sum of terms, each a real number times a whole-number power of $x$.}\\[3pt]
\noindent\textbf{\textcolor{plum}{Degree:}}\\[1pt]
\ansline{The highest exponent on $x$ once the polynomial is written as a sum of terms.}\\[3pt]
\noindent\textbf{\textcolor{plum}{Leading term:}}\\[1pt]
\ansline{The term carrying the highest exponent; it sits first in standard form.}\\[3pt]
\noindent\textbf{\textcolor{plum}{Leading coefficient:}}\\[1pt]
\ansline{The number multiplying the leading term, minus sign included.}\\[3pt]
\noindent\textbf{\textcolor{plum}{Constant term:}}\\[1pt]
\ansline{The term with no $x$ --- it is the value of the polynomial at $x = 0$.}\\[3pt]
\noindent\textbf{\textcolor{plum}{Standard form:}}\\[1pt]
\ansline{Terms written in descending order of exponent: highest power first.}\\[3pt]
\noindent\textbf{\textcolor{plum}{Turning point:}}\\[1pt]
\ansline{A point where the graph stops rising and starts falling, or the reverse.}
\end{vocabbox}

\vspace{0.1in}

\boxguard
\begin{hookbox}
Six expressions are written on the board. Three of them belong to the same family;
three do not.

\begin{center}
$3x^2 - 5$ \qquad $\dfrac{4}{x}$ \qquad $x^3 - 2x + 1$ \qquad
$\sqrt{x} + 6$ \qquad $7$ \qquad $2^x$
\end{center}

Circle the three you think belong together, then write the rule you used.
\ansline{The three circled are $3x^2 - 5$, $x^3 - 2x + 1$, and $7$. In each one every exponent on $x$ is a whole number --- $x$ is never in a denominator, under a radical, or in an exponent.}
\end{hookbox}

\vspace{0.1in}

\boxguard
\begin{notesbox}{1. What Is a Polynomial Function?}
\textbf{Informally.} A polynomial function is what you get when you add up terms,
and every term is built the same way: a \textbf{number} multiplied by $x$ raised to
a \textbf{whole-number} power.

\medskip
For example, $p(x) = 4x^3 - 7x^2 + 2x - 9$ is the sum of four terms:
$4x^3$, \; $-7x^2$, \; $2x$, \; and $-9$.
(The last two are quietly $2x^1$ and $-9x^0$ --- every term really does fit the
pattern.)

\medskip
\textbf{Formally.} A polynomial function can always be written
\[
  p(x) = a_n x^n + a_{n-1}x^{n-1} + \cdots + a_1 x + a_0 .
\]

\begin{itemize}[itemsep=4pt]
  \item Every exponent on $x$ is a \ans{whole} number: $0, 1, 2, 3, \ldots$
  \item Every coefficient $a_n, a_{n-1}, \ldots, a_0$ is a \ans{real} number.
  \item The very first coefficient, $a_n$, is not allowed to be \ans{zero}
        --- otherwise that term would not really be there.
  \item A lone number such as $p(x) = 7$ counts too: it is a polynomial of degree
        \ans{zero}.
\end{itemize}
\end{notesbox}

\vspace{0.1in}

\boxguard
\begin{notesbox}{2. The Parts of a Polynomial, by Name}
Every polynomial has the same four parts, and each one has a name you are expected
to use. Fill in the right column for $p(x) = 4x^3 - 7x^2 + 2x - 9$.

\smallskip
\renewcommand{\arraystretch}{1.5}
\begin{tabular}{|p{0.26\linewidth}|p{0.40\linewidth}|c|}
\hline
\textbf{Part} & \textbf{What it is} & \textbf{For $p(x)$ above} \\
\hline
Degree & the highest exponent on $x$ & \ans{3} \\
\hline
Leading term & the whole term carrying that highest exponent & \ans{$4x^3$} \\
\hline
Leading coefficient & the number in front of the leading term & \ans{4} \\
\hline
Constant term & the term with no $x$ at all & \ans{$-9$} \\
\hline
\end{tabular}
\smallskip

A polynomial is in \textbf{standard form} when the terms run from the
\ans{highest} power down to the \ans{lowest} power. That ordering is worth the
trouble: it puts the two parts we care about most --- the leading term and the
constant term --- at the two ends where they are easy to find.

\medskip
Rewrite each polynomial in standard form, then read off its parts.

\smallskip
\renewcommand{\arraystretch}{1.6}
\begin{tabular}{|p{0.26\linewidth}|p{0.30\linewidth}|c|c|}
\hline
\textbf{As written} & \textbf{Standard form} & \textbf{Degree} & \textbf{Leading coef.} \\
\hline
$5x - 2 + 3x^4 - x^2$ & \ans{$3x^4 - x^2 + 5x - 2$} & \ans{4} & \ans{3} \\
\hline
$9x - x^3 + 2$ & \ans{$-x^3 + 9x + 2$} & \ans{3} & \ans{$-1$} \\
\hline
\end{tabular}
\smallskip

\textbf{Watch the sign.} When a term is written $-x^3$, its coefficient is
\ans{$-1$}, not $1$. The minus sign belongs to the number.
\end{notesbox}

\vspace{0.1in}

\boxguard
\begin{notesbox}{3. What Is \textit{Not} a Polynomial}
This is the skill the rest of the unit leans on. Decide, and say why. The first row
is done for you.

\smallskip
\renewcommand{\arraystretch}{1.6}
\begin{tabular}{|p{0.20\linewidth}|c|p{0.46\linewidth}|}
\hline
\textbf{Expression} & \textbf{Polynomial?} & \textbf{Why / why not} \\
\hline
$5x^3 - x + 8$ & yes & every exponent on $x$ is a whole number \\
\hline
$\dfrac{4}{x}$ & \ans{no} & \ans{negative exponent: $4/x = 4x^{-1}$} \\
\hline
$\sqrt{x} + 6$ & \ans{no} & \ans{fractional exponent: $\sqrt{x} = x^{1/2}$} \\
\hline
$2^x$ & \ans{no} & \ans{the variable sits in the exponent} \\
\hline
$\dfrac{3}{x+1}$ & \ans{no} & \ans{the variable sits in a denominator} \\
\hline
$\dfrac{x^2}{5}$ & \ans{yes} & \ans{dividing by a number is allowed} \\
\hline
$\sqrt{7}\,x^2 - \dfrac{1}{2}x$ & \ans{yes} & \ans{coefficients may be any real numbers} \\
\hline
\end{tabular}
\smallskip

\textbf{The three-question test.} An expression in $x$ is a polynomial only when
all three answers are yes.

\begin{enumerate}[label=(\arabic*), itemsep=4pt]
  \item Is every exponent on $x$ a \ans{whole} number?
        (So no $x^{-2}$ and no $x^{1/2}$.)
  \item Is $x$ kept out of every \ans{denominator}? (So no $\frac{4}{x}$.)
  \item Is $x$ kept out of every \ans{exponent}? (So no $2^x$.)
\end{enumerate}

\medskip
\textbf{The coefficients are never the problem.} $\sqrt{7}$, $-\frac{1}{2}$, and
even $\pi$ are all perfectly legal \ans{coefficients}. Dividing a term by a
\textit{number} is fine too, since $\dfrac{x^2}{5}$ is just $\frac{1}{5}x^2$. Only
the \ans{exponents} on $x$ decide the question.
\end{notesbox}

\vspace{0.1in}

\boxguard
\begin{notesbox}{4. The Family, Sorted by Degree}
The low degrees have names, and you already know most of them.

\smallskip
\renewcommand{\arraystretch}{1.5}
\begin{tabular}{|c|c|l|p{0.32\linewidth}|}
\hline
\textbf{Degree} & \textbf{Name} & \textbf{Example} & \textbf{What its graph looks like} \\
\hline
0 & constant & $p(x) = 6$ & a flat horizontal line \\
\hline
1 & \ans{linear} & $p(x) = 2x - 5$ & a straight slanted line \\
\hline
2 & \ans{quadratic} & $p(x) = x^2 - 4x + 3$ & a parabola \\
\hline
3 & \ans{cubic} & $p(x) = x^3 - 2x$ & new to us this unit \\
\hline
4 & quartic & $p(x) = x^4 - 5x^2 + 4$ & new to us this unit \\
\hline
\end{tabular}
\smallskip

Degrees 0, 1, and 2 are old friends from Algebra 1, Algebra 2, and Unit 1. Lesson
2.1 revisits degree \ans{2} before the unit moves on to degree
\ans{3} and higher.
\end{notesbox}

\vspace{0.1in}

\boxguard[30]
\begin{notesbox}{5. What the Whole Family Looks Like}
No matter the degree, every polynomial graph keeps two habits: it is
\textbf{smooth} (no sharp corners) and \textbf{unbroken} (no holes, no jumps). You
could trace one from left to right without lifting your pencil.

\medskip
A \textbf{turning point} is a spot where the graph stops rising and starts falling,
or stops falling and starts rising.

\begin{center}
\begin{tabular}{ccc}
\begin{tikzpicture}[x=0.50cm, y=0.34cm]
  \draw[->, thick] (-3.2,0) -- (3.2,0);
  \draw[->, thick] (0,-3.2) -- (0,3.4);
  \draw[very thick, plum, domain=-2.8:2.8, samples=80, smooth, variable=\x]
    plot ({\x}, {0.5*\x*\x - 2});
\end{tikzpicture}
&
\begin{tikzpicture}[x=0.50cm, y=0.34cm]
  \draw[->, thick] (-3.2,0) -- (3.2,0);
  \draw[->, thick] (0,-3.2) -- (0,3.4);
  \draw[very thick, plum, domain=-2.6:2.6, samples=80, smooth, variable=\x]
    plot ({\x}, {0.4*\x*\x*\x - 1.6*\x});
\end{tikzpicture}
&
\begin{tikzpicture}[x=0.50cm, y=0.34cm]
  \draw[->, thick] (-3.2,0) -- (3.2,0);
  \draw[->, thick] (0,-3.2) -- (0,3.4);
  \draw[very thick, plum, domain=-2.4:2.4, samples=80, smooth, variable=\x]
    plot ({\x}, {0.35*\x*\x*\x*\x - 1.8*\x*\x + 1});
\end{tikzpicture}
\\[2pt]
\textbf{Graph A} & \textbf{Graph B} & \textbf{Graph C}
\end{tabular}
\end{center}

\smallskip
\renewcommand{\arraystretch}{1.4}
\begin{tabular}{|c|c|c|c|}
\hline
\textbf{Graph} & \textbf{Degree} & \textbf{Turning points you can count} & \textbf{Degree $-\,1$} \\
\hline
A & 2 & \ans{1} & \ans{1} \\
\hline
B & 3 & \ans{2} & \ans{2} \\
\hline
C & 4 & \ans{3} & \ans{3} \\
\hline
\end{tabular}
\smallskip

\textbf{The pattern.} A polynomial of degree $n$ has \textbf{at most}
\ans{$n - 1$} turning points. It may have fewer --- the graph of $p(x) = x^3$ has
none at all --- but it can never have more.

\medskip
Smooth and unbroken is exactly why polynomials are the family calculus starts
with: nothing surprising ever happens between two points on the graph.
\end{notesbox}

\vspace{0.1in}

\boxguard
\begin{notesbox}{6. Where This Unit Is Going}
Today we only named the family. From here: Lesson 2.1 revisits quadratics, the
degree-2 members you already know. Lesson 2.2 measures how fast a polynomial
changes, and 2.3 and 2.4 hunt for the inputs that make a polynomial equal zero
--- first the real ones, then the complex ones. Lesson 2.5 studies what happens at
the far left and far right ends of a graph, 2.6 rewrites polynomials in equivalent
forms, and 2.7 puts it all to work solving equations and inequalities and building
models. Every one of those lessons assumes you can look at an expression and say,
without hesitating, ``that is a polynomial, and its degree is \ldots''
\end{notesbox}

\vspace{0.1in}

\boxguard
\begin{practicebox}
\begin{enumerate}[itemsep=8pt]
  \item Consider $p(x) = -2x^4 + x^3 - 6$.

        \begin{enumerate}[label=(\alph*), itemsep=4pt]
          \item Degree: \ans{4} \quad Leading term: \ans{$-2x^4$} \quad
                Leading coefficient: \ans{$-2$} \quad
                Constant term: \ans{$-6$}
          \item By degree, this polynomial is called a \ans{quartic}, and it has
                at most \ans{3} turning points.
        \end{enumerate}

  \item Is each expression a polynomial? Circle one and give a short reason.

        \begin{enumerate}[label=(\alph*), itemsep=6pt]
          \item $h(x) = 6x^5 - \dfrac{x}{4} + 11$ \quad \ans{Yes} / \textbf{No}

                \vspace{0.05in}
                \ansline{Dividing by the number $4$ is fine: $x/4$ is just $\frac{1}{4}x$.}

          \item $k(x) = 3x^2 + \dfrac{5}{x}$ \quad \textbf{Yes} / \ans{No}

                \vspace{0.05in}
                \ansline{The variable sits in a denominator: $5/x = 5x^{-1}$, a negative exponent.}
        \end{enumerate}
\end{enumerate}
\end{practicebox}

\end{document}
